\chapter*{Abstract}

\addcontentsline{toc}{chapter}{Abstract}

This report presents the work carried out during a six-month final-year internship at Rakuten France, within the IT department's Architecture team, as a Frontend Developer Intern. The internship focused on contributing to the development, modernization, quality, and security of a large-scale frontend platform composed of multiple microfrontend applications.

The first part of the report presents the working environment, including an overview of the company and the Architecture team, followed by a description of the development tools, technologies, and platforms used, as well as the overall architecture of the platform.

The second part presents the main tasks carried out during the internship, namely \textbf{dark mode implementation}, \textbf{modernization of \texttt{next-common} and exploration of Storybook MCP and development tools}, \textbf{development of a visual regression testing infrastructure}, and \textbf{development of a proxy for Google Cloud Storage}. For each task, the report describes its context and motivation, objectives, implementation steps, challenges encountered, and, where relevant, its limitations and potential future improvements. Each task concludes with a brief summary of the work completed and the key lessons learned throughout its implementation.

The report also includes a glossary providing concise definitions of selected technical terms and concepts that may require additional clarification, allowing readers to quickly understand the terminology used throughout the report.

\chapter*{Résumé}
\addcontentsline{toc}{chapter}{Résumé}

Ce rapport présente les travaux réalisés dans le cadre d'un stage de fin d'études de six mois au sein de Rakuten France, dans le département IT et plus précisément au sein de l'équipe Architecture, en tant que stagiaire Développeur Frontend. Ce stage avait pour objectif de contribuer au développement, à la modernisation, à l'amélioration de la qualité et à la sécurisation d'une plateforme frontend à grande échelle, composée de plusieurs applications basées sur une architecture microfrontend.

La première partie du rapport présente l'environnement de travail, notamment une présentation de l'entreprise et de l'équipe Architecture, suivie d'une description des outils de développement, des technologies et des plateformes utilisés, ainsi que de l'architecture globale de la plateforme.

La deuxième partie présente les principales missions réalisées au cours du stage, à savoir \textbf{l'implémentation du mode sombre}, \textbf{la modernisation de \texttt{next-common} et l'exploration de Storybook MCP et de ses outils de développement}, \textbf{la mise en place d'une infrastructure de tests de régression visuelle}, et \textbf{le développement d'un proxy sécurisé pour Google Cloud Storage}. Pour chaque mission, le rapport présente son contexte et ses motivations, ses objectifs, les différentes étapes de sa réalisation, les difficultés rencontrées ainsi que, lorsque cela est pertinent, ses limites et ses perspectives d'amélioration. Chaque mission se conclut par une présentation des principaux enseignements tirés et des résultats obtenus.

Le rapport comprend également un glossaire proposant des définitions synthétiques de certains termes et concepts techniques susceptibles de nécessiter des précisions supplémentaires, afin de permettre au lecteur de comprendre rapidement la terminologie utilisée tout au long du rapport.
